\section{CEP Attributes and Conflict Resolution}
\subsection{Attributes P1 to P7}
\begin{itemize}
\item \textbf{P1 depth of knowledge.} Distributed processing and storage (Spark, Delta, partitioning, broadcast
joins), keyed hashing built from \texttt{sha2}, surveillance statistics (EARS C2, CUSUM), count regression (NB-GLM),
linear programming and visual encoding in D3.
\item \textbf{P2 conflicting requirements.} The five conflicts below, each settled with a number.
\item \textbf{P3 depth of analysis.} Thresholds come from a sweep against injected outbreaks, forecasts from a
52-origin backtest against a baseline, and $\alpha$ from a curve over budgets.
\item \textbf{P4 unfamiliar issues.} Late records (MERGE on the latest report), a facility that goes silent and
back-fills (onset-week counts, missing rather than zero), the same patient at two facilities (DQ-4), and the
Dhaka skew.
\item \textbf{P5 beyond standard codes.} ICD-10 three-character buckets, WHO IDSR~\cite{who2019idsr} for the
diseases and the weekly cadence, and our own suppression and allocation rules where no standard exists.
\item \textbf{P6 stakeholders.} Table~\ref{tab:stakeholders}.
\item \textbf{P7 interdependence.} An error in bronze reaches the dashboard: counting by report date would
create a false outbreak, which would raise forecast demand and move stock. Every gold row carries its bronze
batch IDs, so such an error can be traced back.
\end{itemize}

\subsection{Conflicts C1 to C5}
\textbf{C1 sensitivity against specificity.} We allow at most two false alerts per district-year and take the
most sensitive setting under that limit: EARS C2 at threshold 4.0 finds 97\% of outbreaks with 1.27 false
alerts per district-year\prov{} (Section~\ref{sec:eval}).

\textbf{C2 granularity against privacy.} We publish district by week, not upazila. Every count from 1 to 4 is
hidden, and a second cell is hidden in any row or column with only one hidden cell, so no total reveals it.

\textbf{C3 efficiency against equity.} The LP's equity floor $\alpha$ gives every district a share of each
resource in proportion to its population.
\begin{teammate}{Rafi}
The recommended $\alpha$ and its cost from the curve, once the allocation has real need to allocate.
\end{teammate}

\textbf{C4 freshness against cost.} Forecasts and allocations are weekly decisions, so they run in a daily
batch. Alerts can also run from a Structured Streaming file source that refreshes them as files land
(Section~\ref{sec:data}); between files it holds almost no memory.

\textbf{C5 accuracy against interpretability.} The forecast is a negative-binomial GLM with seasonal terms and
lagged rainfall and temperature: each coefficient can be read as a rate ratio. It beats seasonal naive with a
relative WIS of 0.83\prov.
\begin{teammate}{Rafi}
Confirm whether a more complex model was tried; if not, say so.
\end{teammate}
